% =====================================================================
\section{Introduction}\label{sec:intro}
% =====================================================================
\subsection{Background and Motivation}
A home temperature--humidity sensor answers ``what is the room like now,'' but residents usually need answers to different questions: should the dehumidifier be turned on, is it a good afternoon to dry laundry on the balcony, and will it rain at the destination of a planned outing? The U.S. Environmental Protection Agency recommends keeping indoor RH below 60\%, ideally between 30\% and 50\%~\cite{epa_mold}, yet a single threshold cannot decide an action on its own. Thermal comfort depends on air temperature, mean radiant temperature, humidity, and air speed together with clothing and activity~\cite{ashrae55,mishra2026myths}, and outdoor plans depend on forecasts that the indoor sensor cannot provide.

\subsection{Problem Statement and Research Gap}
Combining these sources is not straightforward. First, the data are \emph{heterogeneous}: an indoor sensor measures one room at a known instant, whereas a township forecast describes a representative point over a validity window~\cite{cwa_township}. Treating the two as the same kind of observation produces recommendations whose basis cannot be checked. Second, operations are \emph{ambiguous}: air conditioning, dehumidification, opening a window, or simply a rising temperature can all lower RH~\cite{nws_dewpoint}, so a curve alone does not reveal what the resident did. Third, \emph{labels are costly}: asking the resident resolves the ambiguity, but frequent questions are intrusive. A general-purpose LLM assistant does not resolve these issues and may state plausible but unsupported values~\cite{ji2023hallucination}.

As summarized in Table~\ref{tab:summary}, existing tool-using LLM agents, smart-home assistants, ambient occupancy models, and thermal comfort tools each address part of the stack, but none jointly validates in-room temperature--humidity sensing, spatiotemporal forecast alignment, time-aligned appliance-operation labels, and evidence-checked recommendations. This paper addresses two questions:
\begin{itemize}
\item \textbf{Q1:} Can an agent select the data that apply to a given room, destination, and time window, and produce recommendations whose every factual element can be traced to a source?
\item \textbf{Q2:} When the indoor environment changes persistently, how can operation labels be obtained at a controlled query burden, and do they improve operation identification over fixed rules?
\end{itemize}

\subsection{Research Contributions}
The main contributions of this work are as follows:
\begin{itemize}
\item \textbf{Evidence-traceable eight-module architecture.} Numerical analysis and rule evaluation are performed by deterministic code, while the LLM is confined to intent parsing, constrained tool routing, and explanation, with a template fallback whenever its output disagrees with the evidence.
\item \textbf{Separated data model and alignment checks.} Indoor measurements, forecast snapshots, user reports, and model inferences are stored separately and checked for quality and spatiotemporal validity so that the agent clarifies or abstains instead of guessing.
\item \textbf{Bounded event-triggered feedback loop.} Event aggregation and a query cooldown record reply time and operation time separately and keep unanswered queries and mixed operations as distinct labels.
\item \textbf{Leakage-free evaluation protocol.} We define a reproducible single-room protocol with event-grouped temporal splits for three comparisons: fixed rules versus a decision tree, two feedback policies, and template versus agent responses. Experimental results are not reported in this draft.
\end{itemize}

\subsection{Paper Organization}
The first version gives advice only and does not control appliances. Section~\ref{sec:related} reviews related work and identifies the research gap. Section~\ref{sec:method} formalizes the problem and presents the proposed eight-module scheme. Section~\ref{sec:experiments} describes the evaluation plan, comparison table, and limitations, and Section~\ref{sec:conclusion} concludes the paper.
